\chapter{Rémunération et fiscalité selon le lieu}\label{ch:in:pay-regulation-and-tax-by-location}

La directive européenne sur les fonds propres de 2013 plafonnait la part variable de la rémunération d'un preneur de risques bancaire au niveau de sa part fixe, ou à deux fois
ce niveau avec l'accord des actionnaires~; le Royaume-Uni a supprimé ce plafond à compter du 31~octobre 2023, et le plafond ne s'applique plus
aux entreprises d'investissement de l'UE depuis 2021. Lorsque le plafond est contraignant, la rémunération glisse du bonus vers le \omterm{def:in:how-pay-works:base}{salaire de base}~; lorsqu'il ne l'est pas, le bonus
peut atteindre plusieurs fois le \omterm{def:in:how-pay-works:base}{salaire de base}, comme l'ont montré les chiffres sur les \omterm{def:in:pay-levels-by-role-firm-type-and-seniority:median}{hauts revenus} du chapitre~14. Et une fois le package fixé, le lieu de résidence du salarié
détermine la part qu'il lui reste~: un package d'un million de dollars laisse environ 516\,000~\$ à Amsterdam et la totalité
à Dubaï, et un régime spécial pour les impatriés peut modifier le résultat de plusieurs dizaines de milliers de dollars. Ce chapitre présente les règles qui
façonnent la rémunération et les impôts qui la prélèvent, pour huit places où le secteur emploie des personnes, et calcule ce que vaut un package
après impôt dans chacune d'elles. Il explique le fonctionnement des systèmes~; il ne constitue pas un conseil fiscal, et la situation personnelle
d'un lecteur dépend de faits que le chapitre ne prend pas en compte.

\section{Régulation des rémunérations~: le plafonnement des bonus, les règles de report et leurs bénéficiaires}

\begin{definition}[Plafonnement des bonus]\label{def:in:pay-regulation-and-tax-by-location:cap}
Un \emph{plafonnement des bonus}\index{plafonnement des bonus} (bonus cap) est une limite légale au rapport entre la part variable et la part fixe de la rémunération. Dans la directive européenne sur les exigences de fonds propres, la
composante variable «~ne doit pas dépasser 100~\% de la composante fixe~» pour chaque personne dont l'activité a un impact significatif sur le
risque de l'établissement, avec une limite pouvant atteindre 200~\% si les propriétaires l'approuvent.
\end{definition}

Trois éléments déterminent si une règle s'applique à une personne donnée.
\begin{itemize}
\item \textbf{Le statut de l'employeur.} Le plafonnement et les règles de report du chapitre~13 s'appliquent aux établissements de crédit et, au Royaume-Uni,
aux banques, aux sociétés de construction et d'épargne ainsi qu'aux entreprises d'investissement désignées par le régulateur~; les entreprises d'investissement de l'UE relèvent
d'un régime propre depuis 2021, sans plafonnement.
\item \textbf{Le rôle de la personne.} Au sein d'une entreprise assujettie, les règles les plus strictes s'appliquent aux preneurs de risques significatifs (livre~16,
chapitre~10), identifiés par leur rôle et leur rémunération.
\item \textbf{Le lieu.} La suppression du plafonnement au Royaume-Uni en 2023 et ses modifications de 2025 concernant les reports s'appliquent aux entreprises régulées au Royaume-Uni~;
le régulateur dont les règles s'appliquent à une personne dépend de l'entité qui l'emploie et du lieu où celle-ci est agréée.
\end{itemize}
Le plafonnement ne limite pas la \omterm{def:in:how-pay-works:base}{rémunération totale}~: une entreprise qui ne peut pas verser un bonus plus élevé peut verser un \omterm{def:in:how-pay-works:base}{salaire de base} plus élevé, ce que les régulateurs britanniques
ont présenté comme l'effet du plafonnement lorsqu'ils l'ont supprimé (livre~16, chapitre~10). Pour un candidat, le plafonnement modifie davantage la structure que le niveau~:
une part plus importante est fixe, ce qui est plus sûr, et une part moindre peut faire l'objet de restitutions ou de reports. Les règles américaines suivent une autre voie~:
les émetteurs cotés doivent récupérer la rémunération incitative après un retraitement comptable (livre~16, chapitre~10).

\begin{example}[Le plafonnement et la structure]\label{ex:in:pay-regulation-and-tax-by-location:mix}
Une banque souhaite proposer à un preneur de risques un total de 1\,000\,000~\euro{} en année normale. Avec le plafond européen sans l'accord des propriétaires, la
part variable ne peut pas dépasser la part fixe, si bien que la part fixe doit être d'au moins 500\,000~\euro{}~; avec l'accord (200~\%), d'au moins
333\,333~\euro{}~; sans plafond, comme au Royaume-Uni depuis 2023, la banque peut proposer, par exemple, 300\,000~\euro{} de part fixe et 700\,000~\euro{} de part
variable. Une année où l'attribution variable est nulle, le salarié perçoit 500\,000~\euro{}, 333\,333~\euro{} ou 300\,000~\euro{}~: le plafond est une assurance
pour le salarié et un coût fixe pour la banque, ce qui explique que la vision qu'en a la banque elle-même, au livre~16,
chapitre~10, porte sur la part de la rémunération qui absorbe les pertes.
\end{example}

\section{Impôt sur le revenu du salaire et du bonus}

\begin{definition}[Taux marginal d'imposition, taux moyen d'imposition]\label{def:in:pay-regulation-and-tax-by-location:rates}
Le \emph{taux marginal d'imposition}\index{taux marginal d'imposition} (marginal tax rate) est la part de la tranche suivante de rémunération prélevée au titre de l'impôt sur le revenu et
des cotisations sociales du salarié~; le \emph{taux moyen d'imposition}\index{taux moyen d'imposition} (average tax rate) est la part de l'ensemble de la rémunération ainsi
prélevée. Un bonus est imposé aux taux marginaux dans lesquels il tombe~; une comparaison de packages utilise le taux moyen.
\end{definition}

Dans tous les systèmes que le chapitre modélise, un bonus est imposé comme un revenu d'emploi, au même titre qu'un salaire~: la retenue à la source sur un bonus peut suivre
des barèmes particuliers au cours de l'année, mais l'impôt de l'année est le même que la rémunération soit versée sous forme de salaire ou de bonus. Les systèmes diffèrent par
leurs taux, leurs seuils, leurs cotisations sociales et par le fait que ces cotisations soient ou non plafonnées.

\begin{dated}{2026-09}{Taux pour un salarié célibataire, 2026}\label{dat:in:pay-regulation-and-tax-by-location:rates}
\footnotesize
\begin{tabular}{@{}p{2.2cm}p{10.4cm}@{}}
\toprule
Londres & 20~\%, 40~\%, 45~\% au-delà de 125\,140~\pounds{}~; abattement personnel de 12\,570~\pounds{}, retiré au-delà de 100\,000~\pounds{}~; cotisations
salariales National Insurance de 8~\% jusqu'à 50\,270~\pounds{}, puis 2~\%.\\
New York & fédéral de 10~\% à 37~\% au-delà de 640\,600~\$~; État de 4~\% à 10.9~\% (9.65~\% à partir de 1\,077\,550~\$), les taux inférieurs
étant repris au-delà de 107\,650~\$ de revenu~; ville de 3.078~\% à 3.876~\%~; Social Security de 6.2~\% jusqu'à 184\,500~\$, Medicare de 1.45~\% et
0.9~\% de plus au-delà de 200\,000~\$.\\
Chicago & fédéral et cotisations sur salaire comme à New York~; Illinois 4.95~\% forfaitaires~; pas d'impôt municipal sur le revenu dans le modèle.\\
Amsterdam & 35.75~\% jusqu'à 38\,883~\euro{}, 37.56~\% jusqu'à 78\,426~\euro{}, 49.5~\% au-delà, cotisations sociales incluses~; crédits
retirés aux revenus plus élevés~; régime expatriés~: jusqu'à 30~\% des salaires exonérés, plafonnés à 262\,000~\euro{} de salaires.\\
Zurich & impôt fédéral jusqu'à 11.5~\% de l'ensemble du revenu au sommet~; impôt cantonal simplifié jusqu'à 13~\%, multiplié par 2.14 (canton 95~\%,
commune 119~\%)~; cotisations sociales de 5.3~\%, et 1.1~\% jusqu'à 148\,200~CHF.\\
Singapour & résidents de 0~\% à 24~\% au-delà de 1\,000\,000~SGD.\\
Hong Kong & le plus faible de 2~\% à 17~\% sur le revenu après abattements et de 15~\% du revenu (16~\% au-delà de 5\,000\,000~HKD)~;
MPF de 5~\% jusqu'à 18\,000~HKD par an.\\
Dubaï & aucun impôt sur le revenu des personnes physiques.\\
\bottomrule
\end{tabular}

\smallskip\noindent Royaume-Uni 2026/27~; fédéral et cotisations sur salaire américaines 2026~; État de New York et ville depuis les instructions 2025~; les
Pays-Bas et Zurich 2026~; Singapour depuis l'année d'imposition 2024~; Hong Kong 2026/27. Chaque paramètre, avec sa source, figure dans
\texttt{data/industry/tax\_2026.csv}.
\end{dated}

La \Cref{fig:in:pay-regulation-and-tax-by-location:avg} montre le taux moyen de 150\,000~\$ à 3~millions de dollars. Les places
se répartissent en trois groupes~: Dubaï à zéro~; Hong Kong et Singapour entre 10~\% et 24~\%, Hong Kong étant presque plat parce que
son taux standard plafonne la moyenne~; et les cinq autres entre 23~\% et 52~\%, en hausse vers leurs taux marginaux les plus élevés.
New York rejoint Londres vers 1~million de dollars, lorsque la tranche de 9.65~\% de l'État et ses reprises entrent en vigueur~; Zurich part bas
et finit près de Chicago.

\begin{omfigure}
\begin{tikzpicture}
\begin{axis}[width=9.2cm,height=6.2cm,xmode=log,xmin=140,xmax=3200,ymin=-2,ymax=56,log ticks with fixed point,
  xtick={150,300,500,1000,2000,3000},xlabel={\small rémunération brute, milliers de \$ (échelle log)},ylabel={\small taux moyen, \%},
  tick label style={font=\footnotesize},grid=major,grid style={gray!20},table/col sep=comma,
  legend style={legend cell align=left,font=\scriptsize,at={(1.02,0.5)},anchor=west}]
\addplot[omDef,thick] table[x=gross,y=new_york]{figdata/industry/15-pay-regulation-and-tax-by-location/avg_rate.csv};
\addplot[omThm,thick] table[x=gross,y=amsterdam]{figdata/industry/15-pay-regulation-and-tax-by-location/avg_rate.csv};
\addplot[black,thick] table[x=gross,y=london]{figdata/industry/15-pay-regulation-and-tax-by-location/avg_rate.csv};
\addplot[omDef,thick,dashed] table[x=gross,y=chicago]{figdata/industry/15-pay-regulation-and-tax-by-location/avg_rate.csv};
\addplot[omThm,thick,dashed] table[x=gross,y=zurich]{figdata/industry/15-pay-regulation-and-tax-by-location/avg_rate.csv};
\addplot[gray,thick] table[x=gross,y=singapore]{figdata/industry/15-pay-regulation-and-tax-by-location/avg_rate.csv};
\addplot[gray,thick,dashed] table[x=gross,y=hong_kong]{figdata/industry/15-pay-regulation-and-tax-by-location/avg_rate.csv};
\addplot[black,thick,dotted] table[x=gross,y=dubai]{figdata/industry/15-pay-regulation-and-tax-by-location/avg_rate.csv};
\addplot[omThm,thick,dotted] table[x=gross,y=amsterdam_expat]{figdata/industry/15-pay-regulation-and-tax-by-location/avg_rate.csv};
\legend{New York,Amsterdam,Londres,Chicago,Zurich,Singapour,Hong Kong,Dubaï,Amsterdam (expatriés)}
\end{axis}
\end{tikzpicture}

\omcaption{Taux moyen de l'impôt sur le revenu et des cotisations salariales sur un package de 150\,000~\$ à 3~millions de dollars, pour un salarié célibataire
sans autre revenu, converti aux taux de change moyens de 2025. Données~: \texttt{firm.aftertax} avec les
paramètres 2026, via \texttt{in\_tax.curve}.}\label{fig:in:pay-regulation-and-tax-by-location:avg}
\end{omfigure}

\section{Régimes des impatriés}

\begin{definition}[Résidence fiscale, régime fiscal des impatriés]\label{def:in:pay-regulation-and-tax-by-location:expat}
La \emph{résidence fiscale}\index{résidence fiscale} est le statut qui rend une personne imposable dans un pays sur ses revenus,
généralement sur les revenus mondiaux, selon les règles du pays sur les jours de présence, le logement principal et les attaches. Un \emph{régime fiscal des
impatriés}\index{régime fiscal des impatriés} est une règle qui impose plus légèrement un résident nouvellement arrivé pendant une durée limitée,
en exonérant une partie du salaire ou des revenus étrangers, afin d'attirer des personnes aux compétences rares.
\end{definition}

Deux régimes illustrent cet éventail.
\begin{itemize}
\item \textbf{Le dispositif néerlandais pour les impatriés} permet à un employeur de verser à un nouvel arrivant éligible jusqu'à 30~\% du salaire en franchise d'impôt pendant au plus
cinq ans, le taux passant à 27~\% à partir du 1\up{er} janvier 2027~; le montant exonéré est plafonné, le taux étant appliqué à
262\,000~\euro{} de salaire au plus, soit 78\,600~\euro{} au maximum. Sur le package d'un million, il fait passer le salaire net de 516\,376~\$ à
560\,341~\$~; son taux moyen passe de 48.4~\% à 44.0~\%, tandis que le taux marginal sur le bonus reste à 49.5~\%, car au-delà du
plafond, chaque euro supplémentaire est imposé en totalité.
\item \textbf{Le Royaume-Uni} a remplacé son régime dit de remittance basis, fondé sur le domicile, par «~un nouveau régime fiscal fondé sur
la résidence à compter du 6 avril 2025~»~: exonération totale «~sur les revenus et plus-values de source étrangère des nouveaux arrivants \dots{} pendant leurs 4 premières années de
\omterm{def:in:pay-regulation-and-tax-by-location:expat}{résidence fiscale}~» après dix ans d'absence. Il allège les revenus et plus-values de l'étranger~; il ne réduit pas à lui seul
l'impôt sur les rémunérations d'un travail exécuté à Londres.
\end{itemize}
D'autres pays ont des régimes qui leur sont propres, assortis de conditions sur les qualifications, la résidence antérieure et la durée~; chacun
est un texte de loi à lire, et non un taux à supposer.

\section{Rémunérations différées et en actions~: quand elles sont imposées}

Aux États-Unis, un bien soumis à un risque substantiel de \omterm{def:in:how-pay-works:rsu}{déchéance} constitue un revenu lors de son vesting, à sa valeur à ce moment
(chapitre~13)~: la rémunération différée est imposée lorsque le salarié la reçoit, et non lorsqu'elle est attribuée. Le modèle du chapitre applique le
même calendrier partout~; d'autres systèmes ont leurs propres règles pour les attributions d'actions, à lire avant tout déménagement. Trois
conséquences en découlent pour un package où la règle s'applique.
\begin{enumerate}
\item \textbf{Le taux est celui de l'année du vesting.} Une attribution accordée une année de faible revenu et dont le vesting intervient une année de
revenus élevés est imposée au taux marginal de cette année-là.
\item \textbf{La valeur est celle du vesting.} Les actions qui montent entre l'attribution et le vesting sont imposées sur la valeur supérieure~;
les actions qui baissent sont imposées sur une valeur moindre, mais le salarié supporte la baisse (le chapitre~10 a montré le cas extrême des jetons).
\item \textbf{Un déménagement entre pays pendant le vesting} peut laisser à plus d'un pays un droit sur la même attribution~;
le chapitre ne modélise pas le partage, qui dépend des règles de chaque pays et des conventions fiscales.
\end{enumerate}

\section{Ce que vaut un montant après impôt}

Le taux marginal est celui auquel est soumis un bonus~; la \cref{fig:in:pay-regulation-and-tax-by-location:marginal} le montre sur toute la
gamme des packages. Il atteint son plafond tôt à Londres (47~\% dès que l'abattement personnel a disparu) et à Amsterdam
(49.5~\%), augmente à New York lorsque la récupération de l'État entre en vigueur (dans le modèle, qui reproduit de façon approchée les tableaux de calcul de l'État,
l'impôt supplémentaire dépasse le salaire supplémentaire sur une tranche de 50\,000~\$ de revenu au-dessus de 1\,077\,550~\$ de revenu imposable), et reste égal ou inférieur à 24~\% à Singapour et à 17~\% à
Hong Kong sur toute la période. Au bas de l'échelle, la réduction progressive des abattements et des crédits pousse le taux marginal au-dessus du
taux légal le plus élevé~: 62~\% à Londres et 56~\% à Amsterdam pour 150\,000~\$.

\begin{center}\footnotesize
\begin{tabular}{@{}>{\raggedright\arraybackslash}p{4.0cm}>{\centering\arraybackslash}p{2.6cm}>{\centering\arraybackslash}p{2.4cm}>{\centering\arraybackslash}p{1.8cm}>{\centering\arraybackslash}p{1.8cm}@{}}
\toprule
Package de 1~million de dollars (30~\% de fixe, 70~\% de bonus) & devise locale & salaire net (\$) & taux moyen & taux marginal\\
\midrule
Dubaï & 1\,000\,000~USD & 1\,000\,000 & 0.0~\% & 0.0~\%\\
Hong Kong & 7\,796\,944~HKD & 844\,104 & 15.6~\% & 16.0~\%\\
Singapour & 1\,305\,833~SGD & 791\,283 & 20.9~\% & 24.0~\%\\
Zurich & 829\,243~CHF & 599\,515 & 40.0~\% & 44.4~\%\\
Chicago & 1\,000\,000~USD & 597\,361 & 40.3~\% & 44.3~\%\\
Londres & 758\,235~GBP & 545\,545 & 45.4~\% & 47.0~\%\\
New York & 1\,000\,000~USD & 540\,584 & 45.9~\% & 50.1~\%\\
Amsterdam & 884\,969~EUR & 516\,376 & 48.4~\% & 49.5~\%\\
\bottomrule
\end{tabular}

\smallskip Impôt sur le revenu et cotisations salariales pour un salarié célibataire sans autre revenu~; taux de change moyens
de 2025~; le taux marginal porte sur les derniers 1\,000~\$ de bonus.
\end{center}

Le salaire net sur le même million va de 516\,376~\$ à Amsterdam à l'intégralité du montant à Dubaï, soit un rapport de 1.94
(la \cref{fig:in:pay-regulation-and-tax-by-location:million}). Parmi les cinq places dont le taux moyen est d'au moins 40~\%,
l'écart est plus étroit~: de 516\,376~\$ à 599\,515~\$, soit 16~\%. La répartition entre fixe et bonus ne change pas ces chiffres, car les deux sont
imposés comme des revenus salariaux~; elle modifie le risque (chapitre~13) et, sous un \omterm{def:in:pay-regulation-and-tax-by-location:cap}{plafonnement des bonus}, ce que l'entreprise peut offrir.

\begin{omfigure}
\begin{tikzpicture}
\begin{axis}[width=9.2cm,height=6.8cm,xmode=log,xmin=140,xmax=3200,ymin=-3,ymax=120,log ticks with fixed point,
  xtick={150,300,500,1000,2000,3000},xlabel={\small salaire brut, en milliers de \$ (échelle log)},
  ylabel={\small taux marginal sur les derniers 5\,000~\$, \%},tick label style={font=\footnotesize},grid=major,
  grid style={gray!20},table/col sep=comma,legend style={legend cell align=left,font=\scriptsize,at={(1.02,0.5)},anchor=west}]
\addplot[omDef,thick] table[x=gross,y=new_york]{figdata/industry/15-pay-regulation-and-tax-by-location/marginal.csv};
\addplot[omThm,thick] table[x=gross,y=amsterdam]{figdata/industry/15-pay-regulation-and-tax-by-location/marginal.csv};
\addplot[black,thick] table[x=gross,y=london]{figdata/industry/15-pay-regulation-and-tax-by-location/marginal.csv};
\addplot[omDef,thick,dashed] table[x=gross,y=chicago]{figdata/industry/15-pay-regulation-and-tax-by-location/marginal.csv};
\addplot[omThm,thick,dashed] table[x=gross,y=zurich]{figdata/industry/15-pay-regulation-and-tax-by-location/marginal.csv};
\addplot[gray,thick] table[x=gross,y=singapore]{figdata/industry/15-pay-regulation-and-tax-by-location/marginal.csv};
\addplot[gray,thick,dashed] table[x=gross,y=hong_kong]{figdata/industry/15-pay-regulation-and-tax-by-location/marginal.csv};
\addplot[black,thick,dotted] table[x=gross,y=dubai]{figdata/industry/15-pay-regulation-and-tax-by-location/marginal.csv};
\legend{New York,Amsterdam,Londres,Chicago,Zurich,Singapour,Hong Kong,Dubaï}
\end{axis}
\end{tikzpicture}

\omcaption{Taux marginal de l'impôt sur le revenu et des cotisations salariales sur les derniers 5\,000~\$ de salaire, par package brut, 2026.
Les pics reflètent de véritables caractéristiques des barèmes~: la suppression de l'abattement personnel au Royaume-Uni entre 100\,000~\pounds{} et
125\,140~\pounds{}, la suppression des crédits néerlandais et les montées en charge de la récupération de New York. Données~: \texttt{in\_tax.marginal\_curve}.}\label{fig:in:pay-regulation-and-tax-by-location:marginal}
\end{omfigure}

\begin{omfigure}
\begin{tikzpicture}
\begin{axis}[width=10.5cm,height=6.4cm,xbar stacked,bar width=9pt,xmin=0,xmax=1080,ymin=0.4,ymax=9.6,ytick=data,
  yticklabels={{Amsterdam},{New York},{Londres},{Amsterdam (régime des expatriés)},{Chicago},{Zurich},{Singapour},{Hong Kong},{Dubaï}},
  yticklabel style={font=\footnotesize},xticklabel style={font=\footnotesize},xlabel={\small en milliers de \$},
  grid=major,grid style={gray!20},table/col sep=comma,legend style={legend cell align=left,font=\scriptsize,at={(0.5,-0.2)},anchor=north,
  /tikz/every even column/.append style={column sep=8pt}},legend columns=3]
\addplot[fill=omDef!55,draw=omDef] table[x=net,y=pos]{figdata/industry/15-pay-regulation-and-tax-by-location/million.csv};
\addplot[fill=omThm!45,draw=omThm] table[x=tax,y=pos]{figdata/industry/15-pay-regulation-and-tax-by-location/million.csv};
\addplot[fill=gray!35,draw=gray] table[x=social,y=pos]{figdata/industry/15-pay-regulation-and-tax-by-location/million.csv};
\legend{salaire net,impôt sur le revenu,cotisations salariales}
\end{axis}
\end{tikzpicture}

\omcaption{Où va un package d'un million de dollars selon le lieu de résidence et de travail, en 2026. Le régime d'impatriés néerlandais
est représenté sur une ligne distincte. Données~: \texttt{in\_tax.million}.}\label{fig:in:pay-regulation-and-tax-by-location:million}
\end{omfigure}

La comparaison se fait en dollars, et les taux de change bougent. Au taux moyen de 2024, une livre sterling s'échangeait contre 1.2785~\$~; au
taux moyen de 2025, 1.3189~\$. Un package londonien libellé en livres valait 3.2~\% de plus en dollars en 2025 qu'en 2024, sans
changement de salaire, et une comparaison établie aux taux d'une année peut s'inverser l'année suivante. Un package est versé dans la monnaie locale~; un
candidat qui compte dépenser ou épargner dans une autre monnaie supporte ce risque de change.

L'impôt n'est pas la seule différence de lieu. Le chapitre~27 confronte ces chiffres au coût de la vie, à la taille
de l'industrie de chaque place et aux trajets domicile-travail~; savoir si un employeur donné dispose d'un desk dans un lieu donné est une question distincte de
ce que l'impôt de la place laisserait.

\section{Tutoriel~: un million, huit façons}\label{tut:in:pay-regulation-and-tax-by-location:tut}

\begin{tutorial}
\textbf{Objectif.} Calculer le salaire net et les taux moyen et marginal de packages allant de 150\,000~\$ à 3~millions de dollars dans huit
lieux. \textbf{État final~:} le tableau daté,
les \cref{fig:in:pay-regulation-and-tax-by-location:avg,fig:in:pay-regulation-and-tax-by-location:million} et le tableau ci-dessus.
\begin{tutsteps}
\item \textbf{Paramètres.} \texttt{data/industry/tax\_2026.csv} contient chaque tranche, taux, cap et crédit, avec un identifiant de registre~;
\texttt{ch\_federal\_tax\_2026.csv} les points de rupture de la table fédérale suisse.
\item \textbf{Structure.} \texttt{firm.aftertax.net\_local(params, place, gross)} applique la structure de chaque système~; à New York, la récupération des
avantages est approchée (\cref{lst:in:pay-regulation-and-tax-by-location:ny}).
\omcode{code/firm/aftertax/firm_aftertax.py}{94}{112}{Impôt de l'État de New York avec la récupération des tranches inférieures, étalée
sur 50\,000~\$, continu à chaque borne de tranche.}\label{lst:in:pay-regulation-and-tax-by-location:ny}
\item \textbf{Convertir et comparer.} \texttt{net\_usd} convertit aux taux moyens de 2025~; \texttt{marginal} mesure
l'écart de salaire net sur les derniers 1\,000~\$.
\end{tutsteps}
\end{tutorial}

\textbf{Pour aller plus loin.} Ajouter la cotisation de retraite professionnelle que Zurich déduit du revenu imposable~; modéliser un
bonus britannique versé au cours de l'année fiscale suivante~; essayer un couple marié à un seul revenu, ce qui modifie la réponse à plusieurs endroits.

\section{Construction~: le calculateur après impôt}\label{bld:in:pay-regulation-and-tax-by-location:aftertax}

\begin{build}
\bfield{Rôle} Ramener sur une base après impôt chaque chiffre de rémunération des chapitres consacrés aux métiers et du choix fait au chapitre~30.
\bfield{Interface} \texttt{firm.aftertax}~: \texttt{load(path, fed\_path)}~; \texttt{progressive(x, brackets)}~;
\texttt{net\_local(params, place, gross, expat)}~; \texttt{local\_per\_usd}~; \texttt{net\_usd}~; \texttt{marginal}~;
\texttt{LOCATIONS}. Bibliothèque standard de Python.
\bfield{Règles} Les paramètres sont des données, chacune assortie d'une source, jamais des littéraux dans le code~; un salarié unique, célibataire, sans autre revenu,
seules les déductions forfaitaires~; bonus et salaire imposés de la même façon~; les simplifications sont nommées dans la documentation du module.
\bfield{Tests d'acceptation} \texttt{code/firm/aftertax/tests/}~: un cas britannique à 100\,000~\pounds et 200\,000~\pounds calculé à la main~; le
calcul de l'impôt fédéral américain sur 1~million de dollars de revenu imposable~; l'impôt de New York continu et plat à 6.85~\% et 9.65~\% là où la récupération
est complète~; l'impôt progressif et l'impôt au taux forfaitaire de Hong Kong égaux au seuil officiel de 2\,132\,500~HKD~; les valeurs de
la table de Singapour~; les valeurs imprimées de la table fédérale suisse, y compris l'arrondi à 100~CHF~; les crédits néerlandais et le cap des impatriés.
\bfield{Pour aller plus loin} Les cas de couples mariés et de familles~; les déductions de retraite~; le calendrier de la retenue à la source du bonus~; un second pays pendant
l'acquisition des droits.
\end{build}

\begin{omsources}
\item HMRC, Rates and thresholds for employers 2026 to 2027; gov.uk, Income Tax rates and Personal Allowances; HM Treasury
et HMRC, Reforming the taxation of non-UK domiciled individuals.
\item IRS, inflation adjustments for 2026 and Topic 751; SSA, contribution and benefit base; New York State, IT-201-I
(2025) and tax expenditure report; Illinois Department of Revenue.
\item Belastingdienst, loonheffingen 2026 tables; Business.gov.nl, the expat scheme.
\item ESTV, Kantonsblatt Zürich and the 2026 federal tax table; Stadt Zürich, Steuerfuss 2026; AHV/IV leaflets 2.01 and
2.08.
\item IRAS, individual income tax rates; GovHK, salaries tax rates; IRD, PAM 61; MPFA; UAE Government portal.
\item Directive 2013/36/EU, article 94; European Banking Authority, high earners 2024.
\end{omsources}

\section{Exercices}

\begin{exercise}[$\star$]\label{exo:in:pay-regulation-and-tax-by-location:1}
Calculer l'impôt sur le revenu britannique et la National Insurance salariale sur un salaire de 100\,000~\pounds en 2026/27.
\end{exercise}

\begin{exercise}[$\star$]\label{exo:in:pay-regulation-and-tax-by-location:2}
Un preneur de risques d'une banque dans l'UE a une rémunération fixe de 400\,000~\euro. Quelle est la rémunération variable maximale que permet le plafonnement, avec et sans
l'approbation des actionnaires~?
\end{exercise}

\begin{exercise}[$\star$]\label{exo:in:pay-regulation-and-tax-by-location:3}
D'après le tableau, quelle place a le taux marginal le plus élevé sur le bonus d'un package de 1~million de dollars, et quel est-il~?
\end{exercise}

\begin{exercise}[$\star\star$]\label{exo:in:pay-regulation-and-tax-by-location:4}
Montrer que l'impôt progressif et l'impôt au taux forfaitaire de Hong Kong sont égaux à 2\,132\,500~HKD, avec l'abattement de base de
145\,000~HKD.
\end{exercise}

\begin{exercise}[$\star\star$]\label{exo:in:pay-regulation-and-tax-by-location:5}
Pourquoi le régime d'impatriés néerlandais abaisse-t-il le taux moyen sur un package de 1~million de dollars, mais pas son taux marginal~?
\end{exercise}

\begin{exercise}[$\star\star$]\label{exo:in:pay-regulation-and-tax-by-location:6}
Pourquoi le taux moyen de New York passe-t-il de 45.9~\% pour 1~million de dollars à 50.1~\% pour 1.5~million de dollars, alors que les taux fédéral et municipal
sont inchangés sur cette plage.
\end{exercise}

\begin{exercise}[$\star\star\star$]\label{exo:in:pay-regulation-and-tax-by-location:7}
\emph{Programmation.} Avec \texttt{firm.aftertax}, trouver, à 1\,000~\$ près, le salaire brut à partir duquel le salaire net de New York
passe pour la première fois sous celui de Londres, et expliquer le croisement.
\end{exercise}

\begin{exercise}[$\star\star\star$]\label{exo:in:pay-regulation-and-tax-by-location:8}
\emph{Trouver la faille.} «~Dubaï paie deux fois plus~: un million y vaut autant que deux millions à Amsterdam.~»
\end{exercise}

\section{Problème~: Un million, huit façons}

\begin{problem}[{Problème du week-end --- un million, huit façons}]\label{pb:in:pay-regulation-and-tax-by-location:1}
Un trader se voit proposer le même package, soit un million de dollars par an avec 30~\% de fixe et 70~\% de bonus, dans huit bureaux
d'une même entreprise. Que laisse chacun, et qu'est-ce qui devrait encore trancher~?

\medskip\noindent\textbf{Partie I --- Les règles.}
\begin{enumerate}
\item Définir le \omterm{def:in:pay-regulation-and-tax-by-location:cap}{plafonnement des bonus} et énoncer celui de l'UE.
\item Qui le plafonnement couvre-t-il et qui en est exclu, depuis 2021 et 2023~?
\item Que fait un plafonnement à la structure de la rémunération et à son niveau~?
\item Définir les taux marginal et moyen d'imposition.
\item La répartition entre fixe et bonus modifie-t-elle l'impôt de l'année dans ces systèmes~? Pourquoi~?
\end{enumerate}

\medskip\noindent\textbf{Partie II --- Les systèmes.}
\begin{enumerate}[resume]
\item Énoncer les taux, l'abattement et la National Insurance du Royaume-Uni pour 2026/27.
\item Énoncer les taux fédéraux, ceux de l'État de New York et ceux de la ville, ainsi que les taxes sur la masse salariale aux États-Unis.
\item Énoncer les taux de la box 1 néerlandaise et les conditions du \omterm{def:in:pay-regulation-and-tax-by-location:expat}{régime fiscal des impatriés}.
\item Énoncer la structure de Zurich~: fédérale, cantonale, municipale, sociale.
\item Énoncer celles de Singapour, de Hong Kong et de Dubaï.
\end{enumerate}

\medskip\noindent\textbf{Partie III --- Les montants.}
\begin{enumerate}[resume]
\item Donner le package dans la devise locale de chaque place, aux taux de 2025.
\item Donner le salaire net, les taux moyen et marginal dans chaque place.
\item Calculer ce qu'ajoute le régime néerlandais des impatriés.
\item Définir la \omterm{def:in:pay-regulation-and-tax-by-location:expat}{résidence fiscale} et un \omterm{def:in:pay-regulation-and-tax-by-location:expat}{régime fiscal des impatriés}~; quelle charge le régime britannique de 2025 allège-t-il~?
\item Quand la rémunération différée est-elle imposée, et quelles trois conséquences en découlent~?
\end{enumerate}

\medskip\noindent\textbf{Partie IV --- Le verdict.}
\begin{enumerate}[resume]
\item Énoncer le \emph{résultat nommé}~: le salaire net tiré du package dans chacune des huit places, et le rapport du plus élevé
au plus faible.
\item Parmi les cinq places dont le taux moyen atteint 40~\% ou plus, quel est le spread~?
\item Qu'omet le modèle qui pourrait modifier le \omterm{def:in:bank-markets-divisions:league}{classement} pour elle~?
\item Pourquoi le \omterm{def:in:bank-markets-divisions:league}{classement} fiscal n'est-il pas celui des lieux où travailler~?
\item En deux phrases, la conseiller.
\end{enumerate}
\end{problem}

\section{Questions d'entretien}

\begin{interviewq}[$\star$ \iqroles{développeur}]\label{iq:in:pay-regulation-and-tax-by-location:1}
Écrivez une fonction de calcul de l'impôt selon un barème progressif donné par tranches. Comment la testez-vous~?
\end{interviewq}

\begin{interviewq}[$\star$ \iqroles{banque, risque}]\label{iq:in:pay-regulation-and-tax-by-location:2}
Pourquoi le Royaume-Uni a-t-il supprimé le \omterm{def:in:pay-regulation-and-tax-by-location:cap}{plafonnement des bonus}, et qu'a conservé l'UE~?
\end{interviewq}

\begin{interviewq}[$\star\star$ \iqroles{chercheur}]\label{iq:in:pay-regulation-and-tax-by-location:3}
La courbe du \omterm{def:in:pay-regulation-and-tax-by-location:rates}{taux moyen d'imposition} s'aplatit dans une place et monte dans une autre. Quelles caractéristiques des deux systèmes expliquent chacune de ces évolutions~?
\end{interviewq}

\begin{interviewq}[$\star\star$ \iqroles{trader}]\label{iq:in:pay-regulation-and-tax-by-location:4}
Votre bonus représente 70~\% de votre rémunération. Préféreriez-vous qu'il soit imposé au taux marginal de l'année d'attribution ou de l'année où il
devient acquis~? Quand la réponse change-t-elle~?
\end{interviewq}

\begin{interviewq}[$\star\star$ \iqroles{développeur, ingénieur ML}]\label{iq:in:pay-regulation-and-tax-by-location:5}
Les paramètres fiscaux changent chaque année. Comment les stockeriez-vous pour pouvoir reproduire exactement le calcul de l'année précédente~?
\end{interviewq}

\begin{interviewq}[$\star\star\star$ \iqroles{chercheur, risque}]\label{iq:in:pay-regulation-and-tax-by-location:6}
Un modèle de rémunération nette d'impôt omet les retraites, les abattements et la situation familiale. Comment borneriez-vous l'erreur pour un seul
\omterm{def:in:pay-levels-by-role-firm-type-and-seniority:median}{haut revenu}, et où est-elle la plus grande~?
\end{interviewq}
